% ---------------------------------------------------------------------
% Das Dokument kompiliert mit pdflatex und ist auf Basis 
% von Koma-Script entstanden. 
%
% Autor des Templates (für Anmerkungen): 
% Michael von Riegen, riegen@informatik.uni-hamburg.de
%
% Einzelne Code-Teile für das Titelblatt sind aus dem Template 
% von Benjamin Kirchheim entnommen.
%
% 25.05.09, Frank Langanke: Vorlage auf aktuelle KOMA-Version aktualisiert
% 26.05.09, Michael von Riegen: Anmerkung --> aktuelles Koma-Script ist nötig!
% 17.10.2016 neues Uni logo
% ---------------------------------------------------------------------
\documentclass[11pt,DIV=15,BCOR=20mm,bibliography=totoc]{scrbook}

% Import von Paketen und Optionen die das gesamte Dokument betreffen
% sind in myPreamble.sty ausgelagert.
\usepackage{myPreamble}
   
% Arbeitet man nur an einem Kapitel, wird durch folgenden Befehl nur dieses eingebunden.
% Spart manuelles auskommentieren von vielen include-Befehlen;
% hat keine Auswirkung auf input-Befehle
% \includeonly{kapitel1}
   
\begin{document}

% TITELSEITE
\begin{titlepage}

	% Fehler "destination with the same identifier" unterdrücken...
  \setcounter{page}{-1}

	% Titelseite
	\begin{figure}[h]
		\begin{minipage}[b]{62mm}
			\includegraphics[width=62mm]{images/unilogo}
		\end{minipage}
		\hspace{4cm}
		\begin{minipage}[b]{59mm}
			\includegraphics[width=59mm]{images/minlogo}
		\end{minipage}
	\end{figure}

	\vfill
	
	\begin{center}
		% Diplomarbeit 
		\noindent { \huge
			Bachelor Thesis \\
		}
		\vspace{14mm}
		% Titel
		\noindent \textbf{\huge
		  Multitouch Robot Control \\
		}
		\vspace{60mm}	
	\end{center}
	
	\vfill
	
	\noindent \textbf{Merlin Steuer} \\
	\noindent \rule{\textwidth}{0.4mm} 
	\noindent{\textrm{2steuer@informatik.uni-hamburg.de}} \\
	\noindent{\textrm{Studiengang Informatik}} \\
	\noindent{\textrm{Matr.-Nr. 6415125}} \\
	\noindent{\textrm{Fachsemester 10}} \\
	\begin{tabbing}
	\hspace{8em} \=  \kill
	Erstgutachter: \> Dr. Norman Hendrich \\
	Zweitgutachter: \> Dennis Kruppke \\
	~ \\
	Abgabe: 23.04.2018
	\end{tabbing}
	
	% Rückseite der Titelseite mit Zitat
	\newpage 
	\thispagestyle{empty}
	\setcounter{page}{0}

	% wenn man Lust auf ein Zitat hat...
	% ... ansonsten auskommentieren
	~\\ \vfill \noindent 
	Mein Dank gilt all denen, die mich hierbei unterstützt haben -- insbesondere Sven, meinen Eltern und Anna-Lia. Danke.
	\textit{-- Merlin Steuer}
\end{titlepage}



% VERZEICHNISSE (Inhaltsverzeichnis, Abkürzungen)
% Vorspann einleiten --> Seitennummerierung römisch
\frontmatter

% Inhaltsverzeichnis
\tableofcontents
\cleardoublepage

% Hauptteil einleiten --> Seitennummerierung wieder arabisch
\mainmatter

\chapter{Introduction}
\section{Motivation}

\section{Objectives of this thesis}

\section{Outline}


\chapter{Basics}

\section{ROS - The Robot Operating System}
\subsection{rosjava / rosandroid}

\section{The Shadow C6 Robotic Hand}
\subsection{Integration into ROS}

\section{Inverse Kinetics with BioIK}

\chapter{Implementation}
\section{Preparations}
\subsection{Setting up ROS}
\subsection{Installing the Android IDE}
\subsection{Compiling rosandroid}

\section{User Interface}

\section{Software Architecture}

\section{Approaches}

\subsection{Grasp Synergies}

\subsection{Direct Finger Positioning}

\chapter{Conclusion}

\section{Outlook}

\cleardoublepage

% VERZEICHNISSE (Abbildungen, Tabellen)
% Literatur 

\bibliographystyle{alphadin}
\bibliography{Bachelorarbeit}
\cleardoublepage

% ERKLÄRUNG
\chapter*{Eidesstattliche Versicherung}
\thispagestyle{empty}
\addcontentsline{toc}{chapter}{Eidesstattliche Versicherung}

Hiermit versichere ich an Eides statt, dass ich die vorliegende Arbeit im
Bachelorstudiengang Informatik selbstständig verfasst und keine anderen als die
angegebenen Hilfsmittel -- insbesondere keine im Quellenverzeichnis nicht
benannten Internet-Quellen -- benutzt habe. Alle Stellen, die wörtlich oder
sinngemäß aus Veröffentlichungen entnommen wurden, sind als solche kenntlich
gemacht.
Ich versichere weiterhin, dass ich die Arbeit vorher nicht in einem anderen Prüfungsverfahren eingereicht habe und die eingereichte schriftliche Fassung der auf dem elektronischen Speichermedium entspricht.

\noindent Ich bin mit einer Einstellung in den Bestand der Bibliothek des Fachbereiches einverstanden.

\vspace{2cm} 

\noindent Reinfeld, den \uline{~~~~~~~~~~~~~~~~~~~~}~~~~~Unterschrift: \uline{~~~~~~~~~~~~~~~~~~~~~~~~~~~~~~~~~~~~~~~~~~~~~~~~~~} 

    
\end{document}